\PassOptionsToPackage{
  unicode,
  pdfusetitle,
  colorlinks,
  linkcolor=vertexDarkRed,
  urlcolor=vertexDarkRed,
  citecolor=vertexDarkRed,
}{hyperref}
\documentclass[
  aspectratio=1610,
]{beamer}
\usetheme{vertex}
\usepackage[american]{babel}

\usepackage[autostyle]{csquotes}

\usepackage{fontspec}
\usepackage{graphicx}

\usepackage[outputdir=build]{minted}
\usepackage{xcolor}
\usepackage{xparse}
\usepackage{tcolorbox}
\tcbuselibrary{minted}
\tcbuselibrary{skins}
\tcbuselibrary{xparse}
\usepackage{accsupp}

\usepackage{bookmark}
\usepackage[shortcuts]{extdash}
\usepackage{fontawesome5}

\newtcblisting{code}[2][]{
  listing engine=minted,
  minted language=#2,
  minted options={
    fontsize=\footnotesize,
    breaklines,
    autogobble,
    linenos,
    numbersep=3mm
  },
  colback=black!5!white,
  colframe=vertexDarkGrey,
  listing only,
  left=5mm,
  enhanced,
  overlay={%
    \begin{tcbclipinterior}
      \fill[black!20!white] (frame.south west) rectangle ([xshift=5mm]frame.north west);
    \end{tcbclipinterior}
  },
  #1
}

\definecolor{positive}{HTML}{15B01A}
\colorlet{negative}{vertexDarkRed}
\setmonofont{Fira Mono}

\newcommand\headline[1]{\textbf{\Large\strut{}#1}\newline}

\DeclareTCBInputListing{inputcode} {O{} m m O{}} {
  listing file=#2,
  listing engine=minted,
  minted language=#3,
  minted options={
    fontsize=\footnotesize,
    breaklines,
    autogobble,
    linenos,
    numbersep=3mm,
    #4
  },
  colback=black!5!white,
  colframe=vertexDarkGrey,
  listing only,
  left=5mm,
  enhanced,
  overlay={%
    \begin{tcbclipinterior}
      \fill[black!20!white] (frame.south west) rectangle ([xshift=5mm]frame.north west);
    \end{tcbclipinterior}
  },
  #1
}

\renewcommand\theFancyVerbLine{%
    \BeginAccSupp{method=escape,ActualText={}}%
    {\ttfamily\tiny\arabic{FancyVerbLine}}%
    \EndAccSupp{}%
}%



\newcommand\headlineframe[1]{%
  \begin{frame}[c]%
    \begin{center}%
      \Huge\color{vertexDarkRed}#1%
    \end{center}%
  \end{frame}%
}%



\author[M. Nöthe]{Maximilian Nöthe}
\title[Testing]{Testing \\ Test Driven Development \\ Continuous Integration}
\date[2021-03-10]{Escape Summer School 2021 – 2021-06-10}
\institute[TU Dortmund]{Astroparticle Physics, TU Dortmund}

\begin{document}
\maketitle

\begin{frame}[c]{overview}
  \tableofcontents
\end{frame}

\section{Introduction}
\headlineframe{Introduction}

\begin{frame}[c]{Automated Software Testing}
  \begin{itemize}
    \item Verifying that a software works as intended is crucial
    \item Doing this manually using whatever method you can think of is very tedious
    \item[$\Rightarrow$] We need automated tests that verify our software
    \item Tests fall into three categories
      \begin{enumerate}
        \item Unit tests
        \item Integration tests
        \item Performance tests
      \end{enumerate}
  \end{itemize}
\end{frame}

\begin{frame}[c]{Unit tests}
  \begin{itemize}
    \item Test single \enquote{units} of the code in isolation
    \item Require modular design of the code base
    \item Are the bedrock of any more complicated tests
    \item Good unit tests serve multiple purposes
      \begin{itemize}
        \item They check correctness of the current code
        \item They provide a safety net for developers changing parts of the code
        \item They provide examples / documentation for how your code is meant to be used
      \end{itemize}
    \item Writing good unit tests can be hard
  \end{itemize}
\end{frame}

\begin{frame}[c]{Frameworks}
  All modern languages have one or more frameworks for tests, a small selection:

  \begin{description}[Python]
    \item[Python] pytest
    \item[C++] Catch2, GoogleTest
    \item[Java] JUnit
    \item[Rust] Part of the language
    \item[Julia] \texttt{Test} module in the standard library
  \end{description}
\end{frame}

\begin{frame}[c]{Integration tests}
  \begin{itemize}
    \item Test that multiple \emph{units} are working together
    \item E.\,g.\ testing a whole command line application
    \item Can grow arbitrarily large / complicated
  \end{itemize}
\end{frame}

\begin{frame}[c]{Performance tests}
  \begin{itemize}
    \item Unit and integration tests usually only test the correctness of code
    \item Performance (regression) tests make sure the code does not get slower
    \item This introduction focuses on unit tests
  \end{itemize}
\end{frame}

\begin{frame}[c, fragile]{Example python code}
  We are going to use this simple function as example for our first unit tests:
  \inputcode[title=\texttt{examples/step1/fibonacci.py}]{./examples/step1/fibonacci.py}{python}
\end{frame}

\section{pytest}
\headlineframe{pytest}

\begin{frame}[c, fragile]{pytest}
  \begin{itemize}
    \item Standard framework for writing unit tests for Python projects
    \item Uses the \mintinline{python}+assert+ statement for tests
    \item Tests fail if an assertion fails or an exception is raised
    \item Uses introspection of the assertion to give detailed error messages
    \item Automatic test detection using patterns:
      \begin{itemize}
        \item Modules matching \mintinline{text}+test_*.py+ or \mintinline{text}+*_test.py+
        \item Functions called \mintinline{text}+test*+
        \item Methods named \mintinline{text}+test*+ of classes named \mintinline{text}+Test*+
      \end{itemize}

    \item Docs: \url{https://pytest.org}
  \end{itemize}

\end{frame}

\begin{frame}[c, fragile]{First Unit Test}
  \inputcode[title=\texttt{examples/step1/test\_fiboncacci.py}]{./examples/step1/test_fibonacci.py}{python}

  \begin{code}{text}
    $ pytest
    =========================== test session starts ==========================
    platform linux -- Python 3.9.4, pytest-6.2.4, py-1.10.0, pluggy-0.13.1
    rootdir: /home/maxnoe/Uni/escape-school2021/testing/slides/examples/step1
    collected 1 item

    test_fibonacci.py .                                               [100%]

    =========================== 1 passed in 0.00s ============================
  \end{code}
\end{frame}

\begin{frame}[c]{Testing Exceptions}
  \inputcode[title=\texttt{examples/step2/fibonacci.py}]{./examples/step2/fibonacci.py}{python}[lastline=4]
  \inputcode[title=\texttt{examples/step2/test\_exception.py}]{./examples/step2/test_exception.py}{python}

  The same can be done for warnings using \mintinline{python}+pytest.warns+
\end{frame}


\begin{frame}[c, fragile]{Careful with floating point numbers}
  \inputcode[title={Naive, this fails}]{./examples/floating/test_floating.py}{python}[firstline=4, lastline=5]
  \inputcode[title={Correct approach, using \texttt{pytest.approx}}]{./examples/floating/test_floating.py}{python}[firstline=7]

  See \url{https://0.30000000000000004.com/}
\end{frame}

\begin{frame}[c, fragile]{Using numpy testing utitlities}
  \inputcode[title={Using numpy}]{./examples/floating/test_numpy.py}{python}
  See \url{https://numpy.org/doc/stable/reference/routines.testing.html}
\end{frame}

\begin{frame}[c, fragile]{Using astropy quantity support}
  \inputcode[title={Using astropy units}]{./examples/floating/test_astropy.py}{python}
\end{frame}

\begin{frame}[c]{Fixtures}
  \begin{itemize}
    \item Data and resources used by tests can be injected into tests using \enquote{fixtures}
    \item Fixtures are provided by functions decorated with \texttt{@fixture}
    \item Fixtures have a scope $⇒$ same object used per session, module, or test
  \end{itemize}

  \inputcode{./examples/fixtures/test_fixtures.py}{python}
\end{frame}


\begin{frame}[c]{Fixtures provided by pytest}
  pytest provides several builtin fixtures for
  \begin{itemize}
    \item temporary directories \mintinline{python}+tmp_path+ / \mintinline{python}+tmp_path_factory+
    \item Testing output to stdout / stderr \mintinline{python}+capsys+
    \item Testing logging \mintinline{python}+caplog+
    \item Monkeypatching \mintinline{python}+monkeypatch+
  \end{itemize}
  More at \url{https://docs.pytest.org/en/6.2.x/fixture.html}
\end{frame}

\begin{frame}[c]{capsys – Fixture for testing the standard streams}
  \inputcode{./examples/fixtures/test_capsys.py}{python}
\end{frame}

\begin{frame}[c]{caplog – Fixture for testing logging}
  \inputcode{./examples/fixtures/test_logging.py}{python}
\end{frame}

\begin{frame}[c]{Temporary paths}
  \begin{itemize}
    \item For tests that need to create files, use the \mintinline{python}+tmp_path+ fixture \\
      $⇒$ Avoids cluttering and conflicts when running tests multiple times / between tests
    \item \mintinline{python}+tmp_path+ has scope test, so each test gets its own temporary directory
    \item These directories are not cleaned up after the tests, so you can inspect the results
    \item If you need a temporary path with a wider scope, add a new fixture using \mintinline{python}+tmp_path_factory+
  \end{itemize}
\end{frame}

\begin{frame}[c]{Temporary paths}
  \inputcode{./examples/fixtures/test_tmp_paths.py}{python}

  Run the test and checkout \mintinline{bash}+/tmp/pytest-of-$USER/pytest-current/test_to_csvcurrent+
\end{frame}


\begin{frame}[c, fragile]{Choosing which tests to run}
  pytest offers fine-grained control over which tests to run

  \begin{itemize}
    \item Select a specific test:
      \begin{code}{bash}
        $ pytest test_module.py::test_name
      \end{code}
    \item Run only tests that failed the last time pytest was run
      \begin{code}{bash}
        $ pytest --last-failed
      \end{code}
  \end{itemize}
\end{frame}

\begin{frame}[c, fragile]{Debugging}
  \begin{itemize}
    \item Unit tests can be very useful for debugging

    \item E.\,g.\ Write a new test that triggers the bug → investigate → make it pass

    \item pytest allows you to jump into pdb when a test fails:
      \begin{code}{bash}
        $ pytest --pdb
      \end{code}

    \item or any other debugger, e.\,g.\ ipython's:
      \begin{code}{bash}
        $ pytest --pdb --pdbcls=IPython.terminal.debugger:TerminalPdb
      \end{code}
  \end{itemize}
\end{frame}

\section{Test Coverage}
\headlineframe{Test Coverage}

\begin{frame}[c, fragile]{Test Coverage}
  \begin{itemize}
    \item Test coverage is a metric measuring how much of the code is tested:
      \begin{equation*}
        \operatorname{coverage} = \frac{\text{Lines of code executed during tests}}{\text{Total lines of code}}
      \end{equation*}
    \item Can be helpful to find parts of code that are not tested (enough).
    \item Especially useful in CI system to check that new / changed code is tested
    \item One more badge \faSmileWink[regular]{}! \includegraphics[height=0.5cm]{badge.pdf}
  \end{itemize}
\end{frame}

\begin{frame}[c, fragile]{Limitations of coverage}
  Hit lines of code are not a perfect measure.
  \begin{code}{python}
    if some_condition is True:
        do_stuff()
    do_other_stuff()
  \end{code}

  \onslide<2->{%
    When during the tests \mintinline{python}+some_condition is True+, this code will have 100\,\% coverage.
    \bigskip
    But what about \mintinline{python}+some_condition is not True+?
  }

  \onslide<3->

  \begin{code}{python}
    result = scipy.optimize.minimize(likelihood, ...)
  \end{code}

  \onslide<4->{Calling functions from other packages can have arbitrarily many branches}
\end{frame}

\section{Mocking}
\headlineframe{Mocking}

\section{Test Driven Development}
\headlineframe{Test Driven Development}

\section{Doctest}
\headlineframe{Doctest}

\section{Continuous Integration}
\headlineframe{Continuous Integration}
\end{document}
